% Cycle de réponse à un appel d'offres
\begin{tikzpicture}[node distance=3.6mm and 8mm]

  \node[btpbox, text width=3.3cm] (n1) {Détection\\[-1pt]\btplabel{\itshape Identification de l'opportunité}};
  \node[btpbox, text width=3.3cm, below=of n1] (n2) {Analyse\\[-1pt]\btplabel{\itshape Lecture et analyse du DAO}};
  \node[btpbox, text width=3.3cm, below=of n2] (n3) {Décision\\[-1pt]\btplabel{\itshape Go / No-go et stratégie}};
  \node[btpbox, text width=3.3cm, below=of n3] (n4) {Mobilisation\\[-1pt]\btplabel{\itshape Constitution de l'équipe}};

  \node[btpboxalt, text width=3.3cm, right=22mm of n1] (n5) {Rédaction\\[-1pt]\btplabel{\itshape Production des pièces de l'offre}};
  \node[btpboxalt, text width=3.3cm, below=of n5] (n6) {Vérification\\[-1pt]\btplabel{\itshape Contrôle de conformité}};
  \node[btpboxalt, text width=3.3cm, below=of n6] (n7) {Soumission\\[-1pt]\btplabel{\itshape Dépôt de l'offre}};
  \node[btpboxalt, text width=3.3cm, below=of n7] (n8) {Suivi\\[-1pt]\btplabel{\itshape Suivi post-soumission}};

  \foreach \i/\j in {1/2, 2/3, 3/4, 5/6, 6/7, 7/8} {
    \draw[btparrow] (n\i) -- (n\j);
  }
  \draw[btparrow] (n4.east) -- ++(0.55,0) |- (n5.west);

  \foreach \n/\num in {n1/1, n2/2, n3/3, n4/4, n5/5, n6/6, n7/7, n8/8} {
    \node[btpstep] at (\n.north west) {\num};
  }

  \begin{scope}[on background layer]
    \node[btpgroup, fit=(n1)(n4), label={[btplabel]below:Phase amont}] {};
    \node[btpgroup, fit=(n5)(n8), label={[btplabel]below:Phase de production}] {};
  \end{scope}

\end{tikzpicture}
